\addtocontents{toc}{\protect\Needspace{3\baselineskip}}
\originalchapter{黑洞、面积界与量子世界面习题}
以下九题取自 8.821 前三组作业. 解答核对官方解答后以中文重新推导, 其中原题省略条件或近似存在问题的地方另加明示. 单位 $\hbar=c=1$, 恢复单位时另说明.

\src{8.821 PS1.3, (a)--(c), 30 分.}
\begin{exercise}
观察者沿径向自由落入 Schwarzschild 黑洞.
(a) 证明到达视界所需固有时有限, Schwarzschild 坐标时无穷.
(b) 证明进入视界后不能向外发送信号, 但仍能接收外来信号.
(c) 证明无法避免到达奇点, 并估计视界后可存活时间. 原题写 $\tau=\pi GM=1.54\times10^{-5}(M/\Ms)\,\mathrm s$, 银河中心 $M\sim4\times10^6\Ms$ 时约一分钟; 官方解答把该量解释为最大值.
\end{exercise}
\begin{solution}
令 $f=1-r_s/r$, $r_s=2GM$, 保持 Killing 能量 $e=f\dot t$. 归一化四速度给 $\dot r^2=e^2-f$. 自外部落入的 $e>0$, 视界附近 $\dot r\to-e$, 固有时增量 $|\dd r|/e$ 有限. 坐标时则 $\dd t/\dd r\sim-1/f$, 积分有对数发散. 这不表示视界有局部物理奇点.

用入射 Eddington--Finkelstein 坐标 $v=t+r_*$, $\dd r_*/\dd r=1/f$, 径向度规 $\dd s^2=-f\dd v^2+2\dd v\dd r$. 入射光 $v$ 不变可以进入, 出射光满足 $\dd r/\dd v=f/2$, 在视界内仍向更小 $r$ 走. 未来指向时向曲线在两条光方向之间, 因而也不能转向外面.

视界内径向自由落体的剩余固有时 $\Delta\tau=\int_0^{r_s}\dd r/\sqrt{e^2-1+r_s/r}\le\int_0^{r_s}\dd r/\sqrt{r_s/r-1}=\pi r_s/2=\pi GM$. 对任意未来时向曲线, 由内部度规的额外空间运动项也得到相同上界. 上界对应理想极限 $e=0$, 自外部一般下落者达不到; 例如从无穷远静止下落 $e=1$ 给 $\Delta\tau=2r_s/3=4GM/3$. 恢复单位上界 $\pi GM/c^3$, 约 $1.55\times10^{-5}(M/\Ms)\,\mathrm s$. 原题等号不能作为所有观察者的精确寿命.
\end{solution}

\src{8.821 PS1.4, (a)--(b), 30 分.}
\begin{exercise}
初始质量 $M$ 的 Schwarzschild 黑洞向真空按完美黑体辐射 $n$ 种非相互作用无质量玻色粒子.
(a) 证明黑洞加辐射的总熵增加.
(b) 估计蒸发寿命, 明确自由度与黑体面积约定.
\end{exercise}
\begin{solution}
发出正能量 $\dd E$ 后黑洞熵减少 $\dd E/T_H$. 辐射有 $s=4u/(3T_H)$, 因而该能量携熵 $4\dd E/(3T_H)$. 总熵变化为 $\dd E/(3T_H)>0$. 不应以过程整体可逆的 $\dd S=\dd E/T$ 代替自由辐射的熵能比.

若每种粒子含一个玻色自由度, 单自由度 Stefan 常数是 $\pi^2/120$. 使用视界面积 $A=16\pi G^2M^2$ 与 $T_H=1/(8\pi GM)$, 功率 $P=n/(30720\pi G^2M^2)$. 积分 $\dot M=-P$, 得 $t_{\rm evap}=10240\pi G^2M^3/n$. 恢复单位后, 寿命为\[t_{\rm evap}=\frac{10240\pi G^2M^3}{n\hbar c^4}.\] 若每种含两偏振, 系数寿命减半. 实际灰体因子、散射有效面积、新粒子阈值和末期量子引力使此数字不精确, 但模型内 $M^3$ 尺度律可复核.
\end{solution}

\src{8.821 PS2.1, 10 分.}
\begin{exercise}
度规 $\dd s^2=-f(r)\dd t^2+\dd r^2/h(r)+\cdots$, $f,h$ 在非极端视界 $r_0$ 有一阶零点, 其余部分正则. 用解析延拓证明 $T_H=\sqrt{f'(r_0)h'(r_0)}/(4\pi)$.
\end{exercise}
\begin{solution}
设外侧 $f'_0,h'_0>0$. $t=-\ii t_E$, 写 $r-r_0=h'_0\rho^2/4$, 径向项成为 $\dd\rho^2$, 时间项成为 $f'_0h'_0\rho^2\dd t_E^2/4$. 角变量 $\vartheta=\sqrt{f'_0h'_0}t_E/2$ 需周期 $2\pi$ 才无锥尖, 故欧氏时间周期 $\beta=4\pi/\sqrt{f'_0h'_0}$, 温度 $1/\beta$ 即题面. 该推导假定时间按选定边界观察者归一化, 对极端高阶零点不可直接取同一局部形式.
\end{solution}

\src{8.821 PS2.4, (a)--(d), 20 分.}
\begin{exercise}
半径 $R$ 球盒内 $Z$ 种自由无质量标量辐射气体.
(a) 用总能 $E$ 表示熵.
(b) 在即将成为半径 $R$ 黑洞的能量下, 求气体与黑洞熵比, 与 Planck 长度比较.
(c) 固定能量密度 $\rho$ 增大系统, 求最大半径 $R_M$, 及该处熵比; $\rho\ll m_P^4$ 时怎样?
(d) 两种情形中取极大的 $Z$ 可形式上超过黑洞熵, 求最小种数的尺度.
\end{exercise}
\begin{solution}
本题 $\rho$ 是\textbf{能量}密度, $l_P=\sqrt G,m_P=1/\sqrt G$. $a_Z=\pi^2Z/30$, $E=a_ZVT^4$, $S=4a_ZVT^3/3$, 消去温度得 $S=\frac43(a_ZV)^{1/4}E^{3/4}$. 坍缩阈值采用 $2GE=R$, $S_{\rm BH}=\pi R^2/G$. 代 $V=4\pi R^3/3$ 得
\[
\frac{S_{\rm gas}}{S_{\rm BH}}=
C Z^{1/4}\sqrt{\frac{l_P}R},\qquad
C=\frac1\pi\left(\frac{64\pi^3}{3645}\right)^{1/4}.
\]
正常有限 $Z$, 宏观 $R\gg l_P$ 时远小于 1. 固定密度时 $2G\rho(4\pi R_M^3/3)=R_M$, 所以 $R_M=\sqrt{3/(8\pi G\rho)}$, 比值为 $CZ^{1/4}(8\pi G^2\rho/3)^{1/4}$, 低于 Planck 密度且种数通常时很小.

形式上要求比值 $>1$, (b) 需 $Z>3645\pi R^2/(64G)$, (c) 需 $Z>10935/(512G^2\rho)$. 原题提示的巨大种数确能在这个自由气体模型中改变比较. 但随物种数改变引力有效截止与黑洞物理, 所以这不是对完整量子引力熵界的构造性反例.
\end{solution}

\src{8.821 PS2.5, (a)--(b), 15 分.}
\begin{exercise}
$SU(N)$ 规范理论闭路径 $C$ 的 Wilson 环 $W_R(C)=\langle\operatorname{Tr}_R\mathcal P\exp(\ii g\oint_C A_\mu\dd x^\mu)\rangle$.
(a) 在大 $N$ 极限求 fundamental 与 adjoint 的关系.
(b) 回忆禁闭理论面积律, 推出两种弦张力关系并解释.
\end{exercise}
\begin{solution}
$F\otimes\bar F=1\oplus{\rm adj}$, 所以对任意输运矩阵 $U$, $\operatorname{Tr}_{\rm adj}U=\operatorname{Tr}_FU\operatorname{Tr}_FU^\dagger-1$. 大 $N$ 因子化给 $W_A=|W_F|^2+O(N^0)$, 领先项为 $N^2$. 若规范化 $w_F=W_F/N,w_A=W_A/(N^2-1)$, 则 $w_A=|w_F|^2+O(N^{-2})$.

在平面禁闭面积律范围 $w_F\sim\exp(-\sigma_F\mathcal A)$, 因而 $\sigma_A\simeq2\sigma_F$, 可理解为两股基本颜色通量. 但是有限 $N$ 的 adjoint 外源可被胶子屏蔽, 真正无限长时间、大距离可能出现断弦与周长行为. 先取 $N\to\infty$ 与先取无限环面积不能无条件互换.
\end{solution}

\src{8.821 PS3.1, 20 分.}
\begin{exercise}
平直目标空间共形规范 Polyakov 作用量具有 Poincaré 全局对称. 找对应的世界面守恒流, 说明物理意义. 原题作用量整体号依其世界面度规约定; 本卷统一用式\eqref{eq:Polyakov}.
\end{exercise}
\begin{solution}
平移 $\delta X^\mu=\epsilon^\mu$, 取 $P^a_\mu=-(2\pi\alpha')^{-1}\partial^aX_\mu$. 波动方程给 $\partial_aP^a_\mu=0$. Lorentz 流是 $J^a_{\mu\nu}=X_\mu P^a_\nu-X_\nu P^a_\mu$, 散度中的梯度乘积反对称组合相消, 剩余场方程项为零. $\int P^\tau\dd\sigma$ 是目标时空动量, $\int J^\tau\dd\sigma$ 是角动量与 boost 荷. 开弦的边界可吸收流; 只有边界条件尊重相应对称时积分荷守恒. 世界面局部应力约束与这些目标空间全局流是不同对象.
\end{solution}

\src{8.821 PS3.2, (a)--(d), 35 分.}
\begin{exercise}
Neumann 开弦用式\eqref{eq:openmodes}的完整模式, $\alpha_0^I=\sqrt{2\alpha'}p^I$, $L_n=\frac12\sum_{I,m}:\alpha_{n-m}^I\alpha_m^I:$, $X^+=2\alpha'p^+\tau$.
(a) 由约束解出 $\alpha_n^-=(L_n-a\delta_{n0})/(\sqrt{2\alpha'}p^+)$, 包含零模, $a=(D-2)/24$.
(b) 用 $[\alpha_m^I,\alpha_n^J]=m\delta_{m+n,0}\delta^{IJ}$ 推导中心荷 $D-2$ 的 Virasoro 代数.
(c) 求 $[L_n,X^I]$, 验证其生成 $\sigma^\pm=\tau\pm\sigma$ 的微分同胚, $\xi_n^\pm=-\ii\ee^{\ii n\sigma^\pm}$.
(d) 解释代数和这些变换的物理意义. 原题 PDF 模展开排版或文本提取中有因子遗漏, 本题采用经约束与正则关系核算的式\eqref{eq:openmodes}.
\end{exercise}
\begin{solution}
(a) $-2\partial_\pm X^+\partial_\pm X^-+(\partial_\pm X^I)^2=0$. 代入模式并比较 Fourier 系数, 横向二次卷积即 $L_n$, 正规序只在零模产生 $a$. 零模给 $p^-=[\alpha'p_I^2+N-a]/(2\alpha'p^+)$, 与 $2p^+p^--p_I^2=(N-a)/\alpha'$ 一致.

(b) 先逐项对易得 $[L_m,\alpha_n^I]=-n\alpha_{m+n}^I$. 因而 $[L_m,L_n]-(m-n)L_{m+n}$ 与所有振子对易, 对非零总模式不能出现常数. 对 $m>0,n=-m$, 在零动量真空计算常数. $L_{-m}|0\rangle$ 的双产生模式为 $\frac12\sum_{k=1}^{m-1}\alpha_{-k}^I\alpha_{-(m-k)}^I|0\rangle$, Wick 收缩给
\[
\langle0|L_mL_{-m}|0\rangle=
\frac{D-2}2\sum_{k=1}^{m-1}k(m-k)
=\frac{D-2}{12}(m^3-m).
\]
最后用 $\sum k=m(m-1)/2$、$\sum k^2=m(m-1)(2m-1)/6$. 这证明所给中心项, 不把它当作经典代数中本已存在的量.

(c) 零模也参与: $[L_n,x^I]=-\ii\sqrt{2\alpha'}\alpha_n^I$, 对 $n=0$ 此式为 $-2\ii\alpha'p^I$. 振子对易连同零模给
$[L_n,X^I]=-\ii(\ee^{\ii n\sigma^+}\partial_++\ee^{\ii n\sigma^-}\partial_-)X^I$.
乘微小参数 $\epsilon$, 等于 $X^I(\sigma^++\epsilon\xi_n^+,\sigma^-+\epsilon\xi_n^-)-X^I$ 的一阶项. 基函数是复变换, 物理实变换由 $n,-n$ 的合适实线性组合给出.

(d) 这些 $L_n$ 是横向自由二维共形场的应力模式, 因而生成剩余共形坐标变换, 中心荷是横向场数 $D-2$. 开弦端点把左右变换约束为同一个组合, 不是闭弦的两个独立 Virasoro 代数. 量子约束和目标 Lorentz 一致性还需临界维数; 此题没有要求独立计算完整 $[M^{-I},M^{-J}]$ 异常.
\end{solution}

\src{8.821 PS3.4, (a)--(d), 30 分.}
\begin{exercise}
玻色弦 $D=26$ 中两 D24 膜原在 $X^{25}=0$, 一膜移至 $a$.
(a) 写连接两膜开弦模式; (b) 求谱.
(c) 检查 $N=1$ 层质量与简并符合 $U(2)$ 世界体作用量 $\mathcal L=\operatorname{Tr}[-F^2/4-(D\phi)^2/2]$, $\langle\phi\rangle=\operatorname{diag}(v,0)$、$v=a/(2\pi\alpha')$ 的 Higgs 相.
(d) 对 $N$ 个 D$p$ 膜在垂直位置 $a_{(k)}^i$, 写标量真空期望值.
\end{exercise}
\begin{solution}
平行方向保留 Neumann 模式. 从 0 到 $a$ 的法向模式为 $X^{25}=a\sigma/\pi+\sqrt{2\alpha'}\sum_{n\ne0}(\alpha_n^{25}/n)\ee^{-\ii n\tau}\sin n\sigma$, 此式实条件 $\alpha_{-n}=(\alpha_n)^\dagger$ 保证实坐标. 相反定向的经典项翻转并加起点 $a$. 伸长贡献给 $M^2=(a/(2\pi\alpha'))^2+(N-1)/\alpha'$.

重合时有四个端点标签, 每个标签 $N=1$ 有 24 个自由度. 分开后两个对角标签仍有 25 维世界体无质量矢量的 23 个偏振及一个标量, 共每个 24. 两个非对角实颜色方向质量为 $|v|$, 每个 massive 矢量有 24 个偏振. 从 $D_\mu\phi=\partial_\mu\phi-\ii[A_\mu,\phi]$, 对角期望值的对易子使非对角 $A$ 获质量 $v^2$, 对应标量被吃掉, 计数与拉伸弦一致. 一般 $\langle\phi^i\rangle=\operatorname{diag}(a_{(1)}^i,\ldots,a_{(N)}^i)/(2\pi\alpha')$, 连接 $k,l$ 的矢量质量为 $|\boldsymbol a_{(k)}-\boldsymbol a_{(l)}|/(2\pi\alpha')$.
\end{solution}

\src{8.821 PS3.5, (a)--(c), 15 分.}
\begin{exercise}
D 膜质量可由开弦真空图求和得到.
(a) 证明弱耦合时膜质量正比 $g_s^{-1}$.
(b) 证明其引力相互作用相对自身质量在弱耦合时可忽略.
(c) 堆叠 $N$ 膜, 何阶 $N$ 使回馈不可忽略?
\end{exercise}
\begin{solution}
每个连通世界面按 Euler 特征数带 $g_s^{-\chi}$, 领先的盘面 $\chi=1$, 所以膜张力正比 $1/g_s$. Newton 常数随 $g_s^2$ 变化, 固定距离和弦尺度下两膜引力能相对单膜质量的比值含 $G T_p\propto g_s$. 堆叠的总张力约 $NT_p$, 无量纲引力回馈参数变成 $g_sN$; $N\sim1/g_s$ 时不再小. $p$、距离、紧致化可给额外几何因子, 此答案是耦合阶数估计. 玻色膜体系含 tachyon, 「近似静止」不能因此升级成严格稳定性结论.
\end{solution}
